\chapter{Prothrombin Time (PT) and INR}

\section*{What Is This Test?}
The prothrombin time (PT) is a coagulation test that measures the time required for plasma to clot after the addition of tissue factor (thromboplastin) and calcium. It primarily evaluates the extrinsic and common pathways of the coagulation cascade: factor VII (extrinsic pathway) and factors X, V, II (prothrombin), and I (fibrinogen) of the common pathway. The PT is most commonly used to monitor warfarin (Coumadin) anticoagulation therapy and to screen for bleeding disorders involving the extrinsic and common coagulation pathways. Because the thromboplastin reagents used in the PT assay vary in sensitivity between manufacturers and lots, the International Normalized Ratio (INR) was developed to standardize PT results across laboratories worldwide. The INR is calculated as: INR = (Patient PT / Mean Normal PT)\textsuperscript{\text{ISI}}, where ISI is the International Sensitivity Index of the thromboplastin reagent. An INR of 1.0 is normal; therapeutic anticoagulation with warfarin typically targets an INR of 2.0--3.0 for most indications (venous thromboembolism, atrial fibrillation) and 2.5--3.5 for mechanical heart valves \cite{keeling2011guidelines}. The PT/INR is also used to assess liver synthetic function (since all coagulation factors except factor VIII are synthesized in the liver) and to evaluate for vitamin K deficiency and disseminated intravascular coagulation (DIC) \cite{tierney2023current}.

\section*{Alternative Names}
PT; Prothrombin Time; INR; International Normalized Ratio; PT/INR; Protime; Quick's Test

\section*{Principle}
The PT is performed by adding thromboplastin (tissue factor, a combination of phospholipid and recombinant or purified tissue factor) and calcium to citrated platelet-poor plasma. Tissue factor binds and activates factor VII, which in turn activates factor X (extrinsic pathway), leading to the conversion of prothrombin (factor II) to thrombin, and ultimately the conversion of fibrinogen to fibrin, producing a clot. The time from addition of thromboplastin-calcium reagent to clot formation is detected optically (increased turbidity) or mechanically (increased impedance to a magnetic stir bar or steel ball). The entire process is performed at 37\textdegree{}C on automated coagulation analyzers (Instrumentation Laboratory ACL TOP, Stago STA-R, Siemens BCS/BCT) \cite{tietz2012textbook}. The normal PT varies by laboratory and reagent but is typically 11--13.5 seconds. The INR standardizes results across different thromboplastin reagents by incorporating the ISI, which is assigned by the manufacturer by comparing their reagent to an international reference thromboplastin (from the World Health Organization). The INR formula: INR = (Patient PT / Geometric Mean Normal PT)\textsuperscript{\text{ISI}}. The geometric mean normal PT is determined by each laboratory from at least 20 healthy normal donors \cite{tripodi2008international}. The ISI is reagent-specific and lot-specific. Sensitive thromboplastins have ISI values close to 1.0, providing more accurate INR results.

\section*{History}
The PT was developed by Dr. Armand Quick at the University of Wisconsin in 1935. Quick was investigating vitamin K and its role in coagulation when he developed a simple test using thromboplastin extracted from rabbit brain as a source of tissue factor. The ``Quick one-stage prothrombin time'' became the standard test for evaluating the extrinsic pathway \cite{quick1935prothrombin}. In 1941, the anticoagulant effect of dicoumarol (a vitamin K antagonist derived from spoiled sweet clover that caused hemorrhagic disease in cattle) was discovered, and warfarin (named after the Wisconsin Alumni Research Foundation) was introduced as a rat poison in 1948 and as a therapeutic anticoagulant in 1954. The PT became the primary method for monitoring warfarin therapy. However, wide variability in thromboplastin reagents led to inconsistent PT results between laboratories. The INR was developed by Dr. Tom Kirkwood under the auspices of the World Health Organization in 1983 \cite{kirkwood1983calibration}. The WHO established international reference thromboplastins, and the ISI concept allowed standardization of warfarin monitoring worldwide. The INR has been universally adopted and is now the standard for reporting warfarin anticoagulation intensity.

\section*{Why This Test Is Ordered}
\begin{itemize}
  \item Monitoring of warfarin (Coumadin) anticoagulation therapy --- regular PT/INR testing every 4--12 weeks once stable, or more frequently during dose adjustments, to maintain therapeutic INR (2.0--3.0 for most indications; 2.5--3.5 for mechanical mitral valve)
  \item Preoperative assessment of coagulation status in patients with suspected bleeding disorder, liver disease, or malnutrition
  \item Evaluation of unexplained bleeding or easy bruising --- prolonged PT with normal aPTT suggests factor VII deficiency or mild vitamin K deficiency
  \item Screening for vitamin K deficiency --- prolonged PT is the earliest coagulation abnormality in vitamin K deficiency because factor VII has the shortest half-life (4--6 hours) of all the vitamin K-dependent factors
  \item Assessment of liver synthetic function --- prolonged PT indicates impaired hepatic synthesis of coagulation factors; used in prognostic scoring of liver disease (MELD score, Child-Pugh score)
  \item Diagnosis of disseminated intravascular coagulation (DIC) --- prolonged PT and aPTT with thrombocytopenia, low fibrinogen, and elevated D-dimer
  \item Detection of congenital factor deficiencies --- factor VII deficiency (isolated prolonged PT), factors II, V, X, or fibrinogen deficiency (prolonged PT and aPTT)
  \item Monitoring of warfarin reversal with vitamin K, fresh frozen plasma, or prothrombin complex concentrate (PCC)
  \item Baseline PT/INR before initiating warfarin therapy
\end{itemize}

\section*{Primary vs Secondary Test}
The PT/INR is a primary coagulation screening test and the standard test for monitoring warfarin therapy and evaluating the extrinsic and common coagulation pathways.

\section*{How This Test Is Performed}
A venous blood sample is collected in a light blue-top tube containing 3.2\% sodium citrate as anticoagulant (9:1 blood-to-citrate ratio) \cite{fischbach2023manual}. The tube must be filled exactly to the marked line to maintain the correct anticoagulant-to-blood ratio. After collection, the tube is gently inverted 3--4 times to mix. The sample is centrifuged at 1,500--2,500 $\times$ g for 15 minutes to obtain platelet-poor plasma (platelet count <10,000/$\mu$L). The plasma is transferred to the automated coagulation analyzer, where it is incubated at 37\textdegree{}C. Thromboplastin (tissue factor-phospholipid reagent) pre-warmed to 37\textdegree{}C and calcium chloride are added simultaneously, triggering the coagulation cascade. The time to clot formation is measured optically or mechanically. The INR is calculated automatically by the analyzer using the geometric mean normal PT and the reagent ISI. The venipuncture procedure takes less than 5 minutes. The analysis is completed within 2--5 minutes after plasma preparation.

\section*{What Preparation Is Needed}
No special preparation is required for the patient. Fasting is not necessary. The patient should inform the healthcare provider of all medications, with particular attention to: warfarin (the indication for the test), heparin (may slightly prolong PT), direct oral anticoagulants (rivaroxaban, apixaban, edoxaban may prolong PT variably depending on the reagent; dabigatran has less effect), antibiotics (broad-spectrum antibiotics reduce vitamin K production by gut flora and may prolong PT), vitamin K supplements (which shorten PT and reduce INR), and drugs that interact with warfarin metabolism (amiodarone, fluconazole, metronidazole, trimethoprim-sulfamethoxazole, NSAIDs, phenytoin, rifampin). Recent intake of large amounts of vitamin K-rich foods (leafy green vegetables, broccoli, Brussels sprouts) should be noted but is not a reason for fasting.

\section*{Contraindications and Precautions}
No absolute contraindications to PT/INR testing. Critical pre-analytical requirements (failure to follow these renders results invalid): the light blue-top (sodium citrate) tube must be filled exactly to the designated line --- the blood-to-citrate ratio must be 9:1; underfilled tubes contain excess citrate relative to blood, chelating all calcium in the sample and causing a falsely prolonged PT (the added calcium in the reagent is insufficient to overcome the excess citrate); overfilled tubes have insufficient citrate, risking clot formation and consumption of coagulation factors; the tube must be gently inverted 3--4 times immediately after collection to ensure adequate mixing; hemolyzed, clotted, or lipemic specimens are unacceptable; the sample must be centrifuged to obtain platelet-poor plasma within 1 hour of collection; if testing cannot be performed within 4 hours, the plasma should be separated and frozen at -20\textdegree{}C or below (PT is stable for 2 weeks frozen); a traumatic venipuncture with tissue factor release activates clotting and shortens PT; samples collected through heparinized lines or catheters are contaminated and produce falsely prolonged PT; patients with polycythemia (hematocrit >55\%) have reduced plasma volume relative to citrate, producing a higher citrate-to-plasma ratio and a falsely prolonged PT --- a special tube with adjusted citrate volume is required; the PT/INR is reagent-dependent and results from different laboratories using different reagents may not be directly interchangeable, emphasizing the importance of the INR for standardization; direct oral anticoagulants (DOACs) variably affect the PT --- rivaroxaban and edoxaban prolong PT in a concentration-dependent manner, but apixaban and dabigatran have minimal effect, and the PT/INR should not be used to monitor DOAC therapy \cite{kitchen2021international}.

\section*{Risks}
Minimal risks of venipuncture: mild pain or stinging, bruising or hematoma at the puncture site, lightheadedness or vasovagal syncope, and rarely, infection or nerve injury. Patients on anticoagulation have an increased risk of hematoma and prolonged bleeding after venipuncture; firm pressure should be applied to the puncture site for 3--5 minutes.

\section*{What Happens After the Test}
PT/INR results are typically available within 30--60 minutes. For warfarin monitoring, the INR result determines the dose adjustment: if INR is below therapeutic range, the weekly warfarin dose is increased by 10--20\%; if INR is above therapeutic range (3.0--4.5 for standard indications), the dose is reduced or held; if INR >4.5 without bleeding, one or more doses are held and the dose is reduced; if INR >10 or serious bleeding is present, vitamin K (oral or IV) and/or prothrombin complex concentrate (PCC) or fresh frozen plasma is administered. For initial evaluation of prolonged PT, the aPTT and platelet count are reviewed. Isolated prolonged PT (normal aPTT) suggests factor VII deficiency or mild vitamin K deficiency. A prolonged PT with prolonged aPTT suggests deficiency of common pathway factors (X, V, II, or fibrinogen), vitamin K deficiency, liver disease, warfarin effect, or DIC. A mixing study (patient plasma mixed 1:1 with normal pooled plasma) distinguishes factor deficiency (correction of PT after mixing) from an inhibitor (failure of correction). Factor assays are then performed to identify the specific factor deficiency.

\section*{Understanding Results}

\subsection*{Below Normal (Shortened PT / INR <1.0)}
\begin{itemize}
  \item \textbf{No clinical significance} in most cases --- a shortened PT is a laboratory finding without known clinical consequences and does not indicate a hypercoagulable state
  \item \textbf{Specimen collection issues:} Traumatic venipuncture with tissue factor release into the sample; partial clot activation in the tube
  \item \textbf{Pregnancy and estrogen therapy:} Elevated factor VII and factor VIII shorten the PT mildly
  \item \textbf{Recent thrombosis or acute illness:} Acute-phase response can shorten the PT
\end{itemize}

\subsection*{Above Normal (Prolonged PT / INR >1.2)}
\begin{itemize}
  \item \textbf{Warfarin (Coumadin) therapy:} Therapeutic anticoagulation produces intentional PT/INR prolongation. Target INR 2.0--3.0 for most indications (venous thromboembolism treatment and prevention, atrial fibrillation, bioprosthetic heart valves); target INR 2.5--3.5 for mechanical mitral valves and some mechanical aortic valves.
  \item \textbf{Vitamin K deficiency:} Prolonged PT is the earliest laboratory abnormality (factor VII half-life is 4--6 hours, the shortest of the vitamin K-dependent factors). Causes include malnutrition, prolonged antibiotic therapy (elimination of vitamin K-producing gut flora), biliary obstruction (impaired fat-soluble vitamin absorption), and malabsorption syndromes (celiac disease, pancreatic insufficiency, short bowel syndrome).
  \item \textbf{Liver disease:} Impaired hepatic synthesis of coagulation factors II, V, VII, IX, X, and fibrinogen (and protein C, protein S, and antithrombin). The PT is a sensitive marker of liver synthetic function and is used in the MELD score (Model for End-Stage Liver Disease) and Child-Pugh classification for cirrhosis prognosis and liver transplant prioritization. Prolonged PT with low factor V (not vitamin K-dependent) helps distinguish liver disease from vitamin K deficiency or warfarin effect.
  \item \textbf{Direct oral anticoagulants:} Rivaroxaban, apixaban, and edoxaban (factor Xa inhibitors) prolong the PT to varying degrees depending on the reagent; dabigatran (direct thrombin inhibitor) minimally affects PT. The PT/INR should not be used to monitor DOAC therapy.
  \item \textbf{Disseminated intravascular coagulation (DIC):} Consumptive coagulopathy with prolonged PT and aPTT, thrombocytopenia, elevated D-dimer, low fibrinogen, and schistocytes on peripheral smear.
  \item \textbf{Super-warfarin poisoning (brodifacoum, bromadiolone):} Long-acting vitamin K antagonist rodenticides cause severe, prolonged PT/INR elevation lasting weeks to months, requiring prolonged high-dose vitamin K therapy.
  \item \textbf{Congenital factor deficiencies:} Factor VII deficiency (isolated prolonged PT, normal aPTT) is the most common congenital extrinsic pathway deficiency; deficiency of factor II (prothrombin), V, X, or fibrinogen prolongs both PT and aPTT.
  \item \textbf{Acquired factor inhibitors:} Rare autoantibodies against specific coagulation factors (acquired hemophilia A with anti-factor VIII, acquired factor V inhibitors).
  \item \textbf{Lupus anticoagulant (rarely):} Some lupus anticoagulants with high antibody titers can prolong the PT if the reagent phospholipid content is low; more commonly, lupus anticoagulants prolong only the aPTT.
  \item \textbf{Massive transfusion:} Dilutional coagulopathy after large-volume packed red blood cell transfusion.
  \item \textbf{Hypothermia:} In vivo or in vitro hypothermia prolongs the enzymatic coagulation reactions.
\end{itemize}

\section*{Normal Results}
\begin{itemize}
  \item \textbf{Prothrombin Time (seconds):} 11.0--13.5 seconds (reagent-dependent; each laboratory establishes its own range for its specific thromboplastin reagent) \cite{wallach2022interpretation}
  \item \textbf{INR (International Normalized Ratio):} 0.9--1.2 (for individuals not on warfarin)
  \item \textbf{Therapeutic INR ranges:}
  \begin{itemize}
    \item Venous thromboembolism treatment and prevention, atrial fibrillation, bioprosthetic heart valve: \textbf{2.0--3.0}
    \item Mechanical mitral valve, mechanical aortic valve with risk factors: \textbf{2.5--3.5}
    \item Mechanical aortic valve (bileaflet, no risk factors): \textbf{2.0--3.0}
  \end{itemize}
  \item \textbf{Newborns and infants:} PT is physiologically prolonged at birth (INR 1.1--1.6) due to low vitamin K-dependent factor levels. PT declines to adult values by 3--6 months. All newborns receive vitamin K injection at birth to prevent hemorrhagic disease of the newborn.
  \item \textbf{Pregnancy (third trimester):} Mildly shortened PT (INR 0.8--1.0) due to increased factor VII.
\end{itemize}

\section*{Follow-up Tests}
\begin{itemize}
  \item \textbf{Activated partial thromboplastin time (aPTT):} To evaluate the intrinsic pathway and determine whether prolongation is limited to the extrinsic pathway (isolated PT prolongation) or involves both pathways (prolonged PT and aPTT)
  \item \textbf{PT mixing study:} Patient plasma mixed 1:1 with normal pooled plasma and PT repeated; correction indicates factor deficiency; failure to correct indicates an inhibitor
  \item \textbf{Individual coagulation factor assays:} Factor VII (isolated PT prolongation); factors II, V, X, and fibrinogen (combined PT and aPTT prolongation); specific factor identified by mixing study pattern
  \item \textbf{Liver function tests (ALT, AST, albumin, bilirubin):} To assess liver synthetic function and diagnose liver disease
  \item \textbf{Serum vitamin K level:} When vitamin K deficiency is suspected
  \item \textbf{Complete blood count with platelet count:} To assess for DIC (thrombocytopenia), liver disease (thrombocytopenia from hypersplenism), or concurrent hematologic disorder
  \item \textbf{Fibrinogen level (Clauss method):} When common pathway deficiency or DIC is suspected
  \item \textbf{D-dimer:} When DIC or thrombotic microangiopathy is suspected
  \item \textbf{Lupus anticoagulant testing (dRVVT, hexagonal phase phospholipid neutralization, silica clotting time):} When a lupus anticoagulant is suspected as cause of prolonged aPTT (usually PT is normal, but rare cases prolong PT)
  \item \textbf{Thrombin time (TT):} To detect heparin contamination, fibrinogen deficiency, dysfibrinogenemia, or fibrin degradation products
\end{itemize}

\section*{References}
\bibliographystyle{unsrtnat}
\bibliography{references}

\clearpage
\section*{Regulatory Status and Authorization Date}
Regulatory status applies to the specific assay, instrument, manufacturer, and intended use rather than to the test name alone. No single approval date can be assigned to this test category. For a commercial assay, verify the current product in the relevant regulator's database: the U.S. Food and Drug Administration (FDA) clearance, approval, or authorization; India's Central Drugs Standard Control Organisation (CDSCO) licence or permission; the European Union In Vitro Diagnostic Regulation (EU IVDR) CE certificate issued through the applicable conformity-assessment route; or the United Kingdom Medicines and Healthcare products Regulatory Agency (MHRA) registration and UKCA/CE status. Laboratory-developed tests may not have an individual FDA, CDSCO, EU IVDR, or MHRA approval date.
\clearpage
